\documentclass[../../../thesis.tex]{subfiles}
\begin{document}
\subsection{Úvodná časť práce}
Hlavným obsahom úvodnej časti sú formálne náležitosti práce a musia byť zaradené v~poradí podľa zoznamu v~úvode tejto kapitoly.

\subsubsection{Obálka, titulný list, zadanie}
Začiatočné stránky práce automaticky generuje univerzitný
informačný systém AIS vo formáte PDF.
Môžeme ich do záverečnej práce vložiť pomocou príkazu \verb|\includepdf| z~balíčka \verb|pdfpages|
alebo využijeme makrá \verb|FEIcover| a~\verb|FEItitle| na vytvorenie obálky a prvej
stránky práce.
Aby boli všetky informácie aktuálne,
treba venovať pozornosť vyplneniu údajových
premenných v~úvode hlavného súboru \verb|thesis.tex|.

Zadanie odporúčame vložiť pomocou spomínaného makra
\verb|\includepdf| tak, že najprv uložíme PDF súbor
so zadaním do priečinka \verb|includes| a prepíšeme
názov súboru v argumente makra.

\subsubsection{Poďakovanie}
Nepovinná, ale veľmi obľúbená časť práce.
Je umiestnené na samostatnej strane zväčša v~dolnej časti.
Jej obsah je ponechaný na autora.
Obsah poďakovania sa nachádza v~súbore
\verb|includes/thanks.tex| a~sadzbu má na starosti
príkaz \verb|\FEIthanks|.

\subsubsection{Slovenský a anglický abstrakt}
Definícia abstraktu vychádza z technickej normy STN ISO 214 Dokumentácia.
Abstrakty (referáty) pre publikácie a dokumentáciu \cite{iso214}.
Termín abstrakt je skrátené, presné vyjadrenie obsahu
bez pridanej interpretácie a kritiky.
Mal by poskytovať čo najviac informácií obsiahnutých v~dokumente.

Abstrakt si netreba zamieňať s termínmi anotácia,
extrakt alebo rezumé.
Anotácia je stručná poznámka, alebo vysvetlenie,
prípadne veľmi stručný opis dokumentu alebo jeho obsahu.
Extrakt predstavuje časti dokumentu vybratých
na reprezentáciu celku.
Rezumé obsahuje stručné zopakovanie významných prínosov
a~záverov v práci.
Nachádza sa zvyčajne na konci dokumentu
a~slúži na doplnenie orientácie čitateľa,
ktorý študoval predchádzajúci text.
Ak je práca napísaná v~anglickom jazyku,
musí obsahovať rezumé v~slovenčine.
V~slovenskej práci nemusí byť rezumé.

\paragraph{Účel a použitie abstraktov}
\begin{itemize}\it
  \item \uv{Dobre vypracovaný abstrakt umožní čitateľom identifikovať
  základný obsah dokumentu, rýchlo a presne stanoviť jeho
  relevanciu, a tak sa rozhodnúť, či potrebujú čítať celý
  dokument.}

  \item \uv{Čitatelia, pre ktorých predstavuje dokument len okrajový
  záujem, často získajú z~abstraktu dostatok informácií a nemusia
  čítať celý dokument.}

  \item \uv{Abstrakty sú často cenné aj pri automatickom vyhľadávaní
  v~plných textoch na získanie predbežných informácií a na
  informačný prieskum.}
\end{itemize}
\begin{flushright}
(Citované z normy STN ISO 214 \cite{iso214})
\end{flushright}

Podľa metodického usmernenia Ministerstva školstva, vedy, výskumu a športu SR č. 56/2011 (čl. 1, ods. 1)
\emph{\uv{abstrakt obsahuje informáciu o cieľoch práce,
jej stručnom obsahu a~v~závere abstraktu
sa charakterizuje splnenie cieľa,
výsledky a~význam celej práce.
Súčasťou abstraktu je 3 -- 5 kľúčových slov.
Abstrakt se píše súvisle ako jeden odsek a jeho rozsah je
spravidla 100 až 500 slov}}~\cite{usmernenie562011}.

Text slovenského a~anglického abstraktu sa nachádzajú
v~súboroch \verb|attachment.tex| a~\verb|attachmentEN.tex| v~priečinku \verb|includes|.

Do dokumentu ich vložia makrá šablóny
\verb|\FEIabstract| a~\verb|\FEIabstractEN|
v~hlavnom súbore \verb|thesis.tex|.
Každé makro má jeden povinný parameter --
cesta a~názov súboru s~textom abstraktu.
Makrá zároveň vytlačia pod abstrakty
v každom jazyku zoznam kľúčových slov.

\subsubsection{Obsah a zoznamy}
Obsah je povinný prehľad jednotlivých kapitol
a~častí práce s~uvedením nadpisov a~strán.
Začína na samostatnej stránke ako nová kapitola
s~nadpisom Obsah bez číslovania,
ktorý sa ale v~samotnom prehľade kapitol nezobrazí.

V~\LaTeX-u zabezpečuje generovanie obsahu príkaz \verb|\tableofcontents|,
ktorý v~mieste použitia vloží automatický zoznam kapitol s~číslami strán.
Obsah vytvára na základe použitia nadpisov \verb|\section|,
\verb|\subsection| a \verb|\subsubsection|.
Toto makro vytvorí v~pracovnom adresári pomocný textový súbor
s~príponou .toc a~na základe neho generuje finálnu podobu obsahu.
Z~toho dôvodu je potrebné kompilátor \LaTeX-u
spustiť minimálne dvakrát za sebou.

V tejto šablóne má na starosti vytvorenie obsahu makro \verb|\FEIcontent|.

\subsubsection*{\normalsize Zoznam ilustrácií,
obrázkov a tabuliek}
Sú to nepovinné prehľady tzv. plávajúcich objektov.
\LaTeX\ pozná na tento účel dva príkazy: \verb|\listoffigures| a~\verb|\listoftables|. Šablóna FEIstyle ponúka alternatívne makrá \verb|\FEIlistOfFigures| a~\verb|\FEIlistOfTables|, ktoré okrem toho nastavia požadovaný typ stránky bez číslovania.

Zvlášť užitočný je príkaz \verb|\FEIlistOfFiguresAndTables|.
Vytvorí totiž spojený zoznam obrázkov a tabuliek s jedným nadpisom.

Ak zoznamy v práci nechceme, môžeme príslušné príkazy z~hlavného súboru \verb|thesis.tex| vymazať alebo ich označiť ako komentár.

\subsubsection*{\normalsize Zoznam skratiek a značiek}
V textových výstupoch vedecko-technických odborov sa používa
množstvo značiek a~skratiek najmä na označenie fyzikálnych
veličín v matematických vzťahoch,
ale aj zostručnenie textového prejavu najmä pri zložitých názvoch
vedeckých metód, zariadení alebo javov.
Sú to napríklad RTG (röntgenové žiarenie),
AFM (mikroskop atómových síl),
TEM (transmisný elektrónový mikroskop),
IR (infračervené žiarenie),
AC (obvod striedavého prúdu) a~mnoho iných.
Ak sa v práci objavia, musí ich autor pri ich prvom výskyte
jasne zadefinovať,
prípadne vysvetliť anglický preklad.
Rovnako to platí pre všetky použité fyzikálne veličiny.

Aj keď je tento zoznam nepovinná súčasť práce,
odporúčame ho zaradiť kvôli lepšej orientácii čitateľa.
Zoznam má podobu slovníka,
značky uvádzame v~abecednom poradí.

Šablóna ponúka dva spôby vytvorenia zoznamu a~práce so skratkami a~značkami v~texte.
\begin{enumerate}
\item Použitie nástrojov balíka \verb|glossary| umožňuje plne automatickú kontrolu nad veľkým množstvom skratiek. Skratky treba najprv definovať v externom súbore \verb|glossary.tex| a~potom ich môžeme v práci používať dvomi spôsobmi. Pri prvom výskyte použijeme skratku aj s~jej opisom, čo zariadi príkaz \verb|acrfull|.
Pri ďalších výskytoch v texte už stačí používať iba skrátený tvar pomocou príkazu \verb|acrshort|.

Po prvom skompilovaní je potrebné spustiť externý program \verb|makeglossaries| a~text skompilovať znova, prípadne kvôli správnemu radeniu strán v~obsahu, treba kompiláciu spustiť aj tretíkrát.
Podobná procedúra sa vyžaduje aj pri práci s~citáciami v systéme Bib\LaTeX.
Makro \verb|\FEIlistOfGlossaries| vytvorí abecedne zoradený zoznam.

Ak sa rozhodneme pre túto možnosť,
treba v hlavnom súbore \verb|thesis.tex| odstrániť
znak komentára pred príkazmi
\verb|\FEIglossaries{includes/glossary}|
a~\verb|\FEIlistOfGlossaries|. Pozor, tieto dva
riadky sa v súbore \verb|thesis.tex|
nachádzajú na rôznych miestach.
Nepremiestňujeme ich.

Balík \verb|glossaries| je nesporne praktická pomôcka, plnohodnotne však funguje iba v anglickom jazyku.
Pri jeho používaní narazíme na problém so skratkami, ktoré pochádzajú z anglických slov.
V slovenskom texte však musíme používať ich terminologické ekvivalenty.
Aj keď si nakoniec vytvoríme slovenský zoznam skratiek, ich automatické použitie bude limitované pri skloňovaní alebo časovaní jednotlivých výrazov.

Viac sa o možnostiach balíka dozvieme z tutoriálu na stránke \href{https://www.ctan.org/pkg/glossaries}{\texttt{www.ctan.org/pkg/ glossaries}}

\item Skratky zadáme manuálne.
Automatické riešenie v predchádzajúcom bode úplne zlyháva pri práci s veličinami, ktorých zoznam predstavuje praktickú pomôcku najmä vo fyzikálnych a~matematických oblastiach techniky. Na označovanie veličín používame rôzne symboly a ich modifikácie, napríklad písmená gréckej abecedy ($\alpha, \omega, \xi$), symboly so šípkami v prípade vektorov ($\vec r, \vec\varphi, \vec\imath$), preškrtnuté h ($\hbar$), zdvojené symboly ako $\mathbb Z$, prípadne aj niečo takéto: $\aleph_0$, čo je hebrejské písmeno alef.

Súbor \verb|manual_glossary.tex| obsahuje príklad, ako by mohol takýto ručne vyrobený zoznam vyzerať.
Makro
\verb|\FEImanualListOfGlossaries|, ktorého parameter je cesta a názov spomínaného súboru, zariadi samotnú sadzbu.
Zoznam si môžeme postupne vytvárať pri písaní a~udržiavať ho v abecednom poradí.
\end{enumerate}

\subsubsection*{\normalsize Zoznamy algoritmov a výpisov kódov programov}
Tieto typy zoznamov vytvoria makrá \verb|\FEIlistOfAlgorithms|, \verb|\FEIlistOfListings| a~sú špecifické pre informatické odbory.

Ak v práci nemáme výpisy kódov alebo algoritmy,
bude potrebné riadky s týmito príkazmi vymazať alebo označiť ako komentár. O~uvádzaní častí kódov
a~zápisov algoritmov píšeme v kapitole \ref{sec:listings}.
\end{document}